\documentclass[review,11pt,3p,authoryear,nopreprintline]{elsarticle}

\usepackage[T1]{fontenc}
\usepackage{amsmath,amssymb,amsthm}
\usepackage{mathrsfs}
\usepackage{bm}
\usepackage{xcolor}
\usepackage[colorlinks=true,citecolor=blue,linkcolor=blue,urlcolor=blue]{hyperref}

\newtheorem{theorem}{Theorem}
\newtheorem{proposition}{Proposition}
\newtheorem{corollary}{Corollary}
\newtheorem{lemma}{Lemma}
\newdefinition{example}{Example}
\newdefinition{definition}{Definition}
\newdefinition{remark}{Remark}
\newdefinition{property}{Property}
\newdefinition{condition}{Condition}
\newtheorem{assumption}{Assumption}

\begin{document}

\begin{frontmatter}
\title{Failure Probability Bounds with Uncertain Surrogate Residual Distributions}

\author[inst1]{Renlong Wang}
\ead{rw106@rice.edu}
\affiliation[inst1]{organization={Department of Civil and Environmental Engineering, Rice University},
  city={Houston},
  postcode={TX 77005},
  country={USA}}

\begin{abstract}
\end{abstract}

\begin{keyword}
\end{keyword}
\end{frontmatter}

\section{Introduction}

\section{Problem Description}
\label{sec-problem}

For a random input $X\sim\mathbb{P}$ and a measurable limit-state function $g:\mathcal{X}\to\mathbb{R}$, the failure probability is defined as
\[
  p_f
  := \Pr\{g(X)\leq 0\}
  = \mathbb{E}_{X\sim\mathbb{P}} \left[\mathbf{1}_{\{g(X)\leq 0\}}\right]
  = \int_{\mathcal{X}}
    \mathbf{1}_{\{g(x)\leq 0\}}\,\mathbb{P}(\mathrm{d}x).
\]
For fixed $g$, this functional is linear in $\mathbb{P}$.

Evaluating $g$ can be computationally expensive, particularly when each evaluation requires a high-fidelity simulation. A common approach is to replace $g$ with a computationally
efficient surrogate model. Let
\[
  \mathcal{D}
  := \{(x_i,g(x_i))\}_{i=1}^{m}
\]
denote the training dataset used to construct the surrogate $\hat{g}_{\mathcal{D}}$. Throughout the following analysis, $\mathcal{D}$ and $\hat{g}_{\mathcal{D}}$ are treated as fixed. For every $x\in\mathcal{X}$, the residual function is defined as
\[
  r_{\mathcal{D}}(x)
  := g(x)-\hat{g}_{\mathcal{D}}(x) .
\]

Although $r_{\mathcal{D}}$ is deterministic for fixed $g$ and $\mathcal{D}$, its values are generally unknown at unevaluated inputs. In surrogate-based reliability analysis, uncertainty about the residual function is often represented by a real-valued stochastic process
\[
  R_{\mathcal{D}}
  := \{R_{\mathcal{D}}(x):x\in\mathcal{X}\},
\]
whose sample paths represent possible residual functions consistent with the observed training data. Let $\mathcal{R}$ be a measurable space of admissible residual functions such that the evaluation map $(x,r)\mapsto r(x)$ is measurable. We denote the probability law of $R_{\mathcal{D}}$ on $\mathcal{R}$ by $\mathbb{Q}$. All process laws are understood conditional on the fixed training dataset $\mathcal{D}$.

Conditional on the fixed training dataset $\mathcal{D}$, the input $X\sim\mathbb{P}$ is assumed to be independent of the residual process $R_{\mathcal{D}}$. The corresponding surrogate-based failure probability is
\begin{equation} 
    p_{\hat{f}}(\mathbb{Q}) 
    := \Pr_{\mathbb{P}\otimes\mathbb{Q}}
       \left\{
         \hat{g}_{\mathcal{D}}(X)+R_{\mathcal{D}}(X)\leq 0
       \right\} 
    = \int_{\mathcal{R}}
         \int_{\mathcal{X}}
         \mathbf{1}_{\{
           \hat{g}_{\mathcal{D}}(x)+r(x)\leq 0
         \}}
         \mathbb{P}(\mathrm{d}x)
       \mathbb{Q}(\mathrm{d}r). 
\end{equation}
Thus, $p_{\hat{f}}(\mathbb{Q})$ averages the failure probability over possible residual functions. For fixed $\mathbb{P}$ and $\hat{g}_{\mathcal{D}}$, it is linear in $\mathbb{Q}$.

However, in complex engineering systems, the available information may be insufficient to specify the probability law of the residual process reliably. We therefore construct a nonempty ambiguity set $\mathscr{Q}$ of candidate process laws consistent with the available information. Each $\mathbb{Q}\in\mathscr{Q}$ represents a plausible description of uncertainty about the residual function.

For fixed $\mathbb{P}$ and $\hat{g}_{\mathcal{D}}$, we seek the range of surrogate-based failure probabilities over all candidate process laws in $\mathscr{Q}$. The upper
and lower bounds are
\begin{equation}
  \overline{p}_{\hat{f}}
  := \sup_{\mathbb{Q}\in\mathscr{Q}}
     p_{\hat{f}}(\mathbb{Q}),
  \label{eq-upper-bound}
\end{equation}
and
\begin{equation}
  \underline{p}_{\hat{f}}
  := \inf_{\mathbb{Q}\in\mathscr{Q}}
     p_{\hat{f}}(\mathbb{Q}),
  \label{eq-lower-bound}
\end{equation}
respectively. These bounds characterize the range of surrogate-based failure probabilities over the candidate process laws in $\mathscr{Q}$. Their computation generally requires optimization over probability measures on a function space, leading to infinite-dimensional optimization problems.

Accordingly, success depends on two components:
\begin{enumerate}
  \item \textbf{Statistically calibrated ambiguity sets.} Select suitable residual distribution features and calibrate the set size using confidence regions to obtain finite-sample guarantees.
  \item \textbf{Tractable reformulations.} Reformulate both bound problems for efficient computation without additional approximation error.
\end{enumerate}

\section{Ambiguity Set Design}

Gaussian process models induce multivariate Gaussian
distributions for residuals evaluated at finitely many inputs.
When information about the joint residual distribution is
limited, estimates of its mean and covariance matrix provide
a practical basis for characterizing uncertainty.
We construct a moment-based ambiguity set to account for
uncertainty in both the distributional form and the estimated
moments.

The admissible residual family $\mathcal{R}$ is chosen to be
consistent with the training data: every $r\in\mathcal{R}$
satisfies
\[
  r(x_i^{\mathrm{tr}})
  =
  g(x_i^{\mathrm{tr}})
  -\hat{g}_{\mathcal{D}}(x_i^{\mathrm{tr}}),
  \qquad i=1,\ldots,m.
\]
Choose reconstruction nodes $x_1,\ldots,x_n\in\mathcal{X}$
and define
\[
  \epsilon
  := \left(
    R_{\mathcal{D}}(x_1),\ldots,R_{\mathcal{D}}(x_n)
  \right)^{\top}
  \in\mathbb{R}^{n}.
\]
The training inputs impose the observed residual values,
while the reconstruction nodes identify the remaining
variation within $\mathcal{R}$.
Let $\Xi\subseteq\mathbb{R}^{n}$ be a known nonempty Borel
set of admissible nodal residual vectors.

\begin{assumption}[Finite-dimensional residual representation]
\label{ass-residual-reconstruction}
  The evaluation map
  \[
    r\longmapsto
    \left(r(x_1),\ldots,r(x_n)\right)^{\top}
  \]
  is a measurable bijection from $\mathcal{R}$ onto $\Xi$
  with a measurable inverse $S:\Xi\to\mathcal{R}$.
\end{assumption}

For an affine residual family with $n$ free parameters,
this condition follows when the nodal evaluation matrix
is invertible.
The map $S$ reconstructs an admissible residual function
from its nodal values.

Define $h:\Xi\to[0,1]$ by
\begin{equation}
  h(\epsilon)
  :=
  \int_{\mathcal{X}}
    \mathbf{1}_{\{
      \hat{g}_{\mathcal{D}}(x)+S(\epsilon)(x)\leq0
    \}}\,
    \mathbb{P}(\mathrm{d}x).
  \label{eq-residual-failure-function}
\end{equation}
Thus, reliability is evaluated under $\mathbb{P}$ using
the residual function reconstructed from $\epsilon$.
For a process law $\mathbb{Q}$, write
$\mathbb{Q}_{\epsilon}
:=\operatorname{Law}_{\mathbb{Q}}(\epsilon)$
for the induced joint distribution.

\begin{definition}[Moment-based ambiguity set]
  Let $\mathcal{F}$ be the Borel $\sigma$-algebra on
  $\mathbb{R}^{n}$, and let $\mathscr{M}$ denote the set
  of probability measures on $(\mathbb{R}^{n},\mathcal{F})$
  with finite second moments.
  Given a mean vector estimate $\mu_0$, a positive definite
  covariance matrix estimate $\mathbf{\Sigma}_0$, and
  parameters $\gamma_1\geq0$ and $\gamma_2\geq1$, define
  \begin{equation}
    \mathscr{Q}_{\epsilon}
    := \left\{
      \mathbb{Q}_{\epsilon}\in\mathscr{M}
      \;\middle|\;
      \begin{aligned}
        & \left(
            \mathbb{E}_{\mathbb{Q}_{\epsilon}}[\epsilon]-\mu_0
          \right)^{\top}
          \mathbf{\Sigma}_0^{-1}
          \left(
            \mathbb{E}_{\mathbb{Q}_{\epsilon}}[\epsilon]-\mu_0
          \right)
          \leq\gamma_1, \\
        & \mathbb{E}_{\mathbb{Q}_{\epsilon}}\!\left[
            (\epsilon-\mu_0)(\epsilon-\mu_0)^{\top}
          \right]
          \preceq\gamma_2\mathbf{\Sigma}_0, \\
        & \mathbb{Q}_{\epsilon}(\Xi)=1
      \end{aligned}
    \right\}.
    \label{eq-moment-ambiguity-set}
  \end{equation}
  The parameters are chosen so that
  $\mathscr{Q}_{\epsilon}$ is nonempty.
\end{definition}

The first constraint controls uncertainty in the residual
mean vector. The second bounds the second-moment matrix
about $\mu_0$. The third specifies the residual support.

We define the corresponding process-law ambiguity set
$\mathscr{Q}$ to contain all probability laws $\mathbb{Q}$
on $\mathcal{R}$ whose induced joint distributions satisfy
$\operatorname{Law}_{\mathbb{Q}}(\epsilon)
\in\mathscr{Q}_{\epsilon}$.

\begin{proposition}[Exact finite-dimensional representation]
\label{prop-exact-finite-representation}
  Under Assumption~\ref{ass-residual-reconstruction},
  \begin{equation}
    \mathscr{Q}_{\epsilon}
    =
    \left\{
      \operatorname{Law}_{\mathbb{Q}}(\epsilon):
      \mathbb{Q}\in\mathscr{Q}
    \right\},
    \label{eq-exact-induced-ambiguity-set}
  \end{equation}
  and the failure probability bounds satisfy
  \begin{equation}
    \begin{aligned}
      \overline{p}_{\hat{f}}
      &=
      \sup_{\mathbb{Q}_{\epsilon}\in\mathscr{Q}_{\epsilon}}
      \mathbb{E}_{\mathbb{Q}_{\epsilon}}[h(\epsilon)], \\
      \underline{p}_{\hat{f}}
      &=
      \inf_{\mathbb{Q}_{\epsilon}\in\mathscr{Q}_{\epsilon}}
      \mathbb{E}_{\mathbb{Q}_{\epsilon}}[h(\epsilon)].
    \end{aligned}
    \label{eq-finite-dimensional-bounds}
  \end{equation}
\end{proposition}

\begin{proof}
  By Assumption~\ref{ass-residual-reconstruction}, each
  admissible residual function is reconstructed from its
  nodal values. Hence,
  \[
    p_{\hat{f}}(\mathbb{Q})
    =
    \mathbb{E}_{\mathbb{Q}_{\epsilon}}[h(\epsilon)].
  \]
  Every $\mathbb{Q}\in\mathscr{Q}$ induces a distribution
  in $\mathscr{Q}_{\epsilon}$ by construction.
  Conversely, for any
  $\mathbb{Q}_{\epsilon}\in\mathscr{Q}_{\epsilon}$,
  take $\epsilon\sim\mathbb{Q}_{\epsilon}$ and set
  $R_{\mathcal{D}}=S(\epsilon)$.
  Its evaluation vector has distribution
  $\mathbb{Q}_{\epsilon}$, and its process law belongs
  to $\mathscr{Q}$.
  This proves the equality of the induced distribution
  sets and the stated representations of both bounds.
\end{proof}



% \section{Ambiguity Set Design}

% Gaussian distributions have been used to model surrogate residuals in reliability analysis. When information about the residual distribution is limited, its estimated mean and covariance provide a practical basis for characterizing uncertainty. We therefore construct the following moment-based ambiguity set to account for uncertainty in both the distributional form and the estimated moments.

% \begin{definition}[Moment-based ambiguity set] 
%   Let $\mathcal{F}$ its Borel $\sigma$-algebra and $\mathscr{M}$ the set of probability measures on $(\mathbb{R},\mathcal{F})$ with finite second moments. Given estimates $\mu_0$ and $\mathbf{\Sigma}_0$ of the residual mean and covariance matrix, respectively, and parameters $\gamma_1\geq0$ and $\gamma_2\geq1$, define
%   \begin{equation}
%     \mathscr{Q} := \left\{
%       \mathbb{Q} \in \mathscr{M}
%       \left|
%       \begin{aligned}
%         & \left(\mathbb{E}_{\mathbb{Q}}[\epsilon]-\mu_0\right)^{\top} \mathbf{\Sigma}_0^{-1} \left(\mathbb{E}_{\mathbb{Q}}[\epsilon]-\mu_0\right) \leq \gamma_1 \\
%         & \mathbb{E}_{\mathbb{Q}} \left[ (\epsilon-\mu_0)(\epsilon-\mu_0)^{\top} \right] \preceq \gamma_2\mathbf{\Sigma}_0 \\
%         & \Pr \left\{\epsilon \in \Xi\right\} = 1
%       \end{aligned}
%       \right.
%     \right\}.
%   \end{equation} 
% \end{definition}

% The first constraint in the moment-based ambiguity set assumes that the mean of $\epsilon$ lies in an ellipsoid of size $\gamma_1$ centered at the estimate $\mu_0$. The second constraint bounds the second-moment matrix about $\mu_0$ from above by $\gamma_2\mathbf{\Sigma}_0$ in the positive semidefinite order, thereby controlling both residual variability and mean deviation from $\mu_0$ in every direction. The parameters $\gamma_1$ and $\gamma_2$ control the allowable mean deviation and the second-moment bound about $\mu_0$, respectively.

% Before calibrating the size parameters of the moment-based ambiguity set, we state the following finite-sample guarantee, which does not relies on the specify form of ambiguity set.
% \begin{theorem} 
%   For $\delta\in(0,1)$, if the ambiguity set $\mathscr{Q}$ contains the true distribution $\mathbb{Q}^{\star}$ with probability at least $1 - \delta$, i.e.,
%   \[
%     \Pr \left\{ \mathbb{Q}^{\star} \in \mathscr{Q} \right\} \geq 1 - \delta,
%   \]
%   then
%   \[
%     \Pr \left\{ \underline{p}_{\hat{f}} \leq p_{\hat{f}}(\mathbb{Q}^{\star}) \leq \overline{p}_{\hat{f}} \right\} \geq 1 - \delta.
%   \]
% \end{theorem}

% \begin{proof}
%   By the definition of the lower and upper bounds, $\mathbb{Q}^{\star} \in \mathscr{Q}$ imples
%   \[
%     \underline{p}_{\hat{f}} \leq p_{\hat{f}}(\mathbb{Q}^{\star}) \leq \overline{p}_{\hat{f}}.
%   \]
%   Therefore,
%   \[
%     \Pr\left\{ \underline{p}_{\hat{f}} \leq p_{\hat{f}}(\mathbb{Q}^{\star}) \leq \overline{p}_{\hat{f}} \right\}
%     \geq
%     \Pr\left\{\mathbb{Q}^{\star}\in\mathscr{Q}\right\} \geq 1-\delta.
%   \]
% \end{proof}

% The following proposition states conditions under which the true probability distribution of $\epsilon$ falls in the ambiguity set $\mathscr{Q}$ by a proper setting of $\gamma_1$ and $\gamma_2$ for the fixed sample size $N$ and confidence level $1 - \delta$.

% \begin{proposition}
%   For a given $\delta$, let 
%   \begin{equation}
%     \gamma_1^{N} = \qquad \gamma_2^{N} = . 
%   \end{equation}
%   There then exists a positive integer $N_0$ depending on $\delta$ such that, for the true probability distribution $\mathbb{Q}^{\star}$ of surrogate residual $\epsilon$,
%   \[
%     \Pr\left\{\mathbb{Q}^{\star}\in\mathscr{Q}(\gamma_1^{N},\gamma_2^{N})\right\} \geq 1-\delta
%   \]
%   for all $N \geq N_0$.
% \end{proposition}

% \section{Tractable Reformulation}

\end{document}
